\documentclass[a4paper,12pt]{article}
\usepackage[T2A]{fontenc}
\usepackage[utf8]{inputenc}
\usepackage[russian]{babel}
\usepackage{amsmath,amssymb,mathtools}
\usepackage{helvet}
\renewcommand{\familydefault}{\sfdefault}
\usepackage{icomma}
\usepackage[a4paper,top=23mm,bottom=24mm,left=23mm,right=23mm,headheight=25pt]{geometry}
\usepackage{microtype}
\usepackage{needspace}
\usepackage{tikz}
\usetikzlibrary{arrows.meta,positioning,shapes.geometric,calc}
\usepackage{booktabs,tabularx,array}
\usepackage{enumitem}
\usepackage{fancyhdr}
\usepackage{setspace}
\usepackage{xcolor}
\usepackage[hidelinks]{hyperref}
\definecolor{accent}{HTML}{008080}
\definecolor{darkaccent}{HTML}{005B61}
\definecolor{textgray}{HTML}{333333}
\definecolor{muted}{HTML}{636D70}
\setlength{\parindent}{0pt}
\setlength{\parskip}{5.5pt plus 1pt minus 1pt}
\onehalfspacing
\setlist[itemize]{leftmargin=6mm,topsep=2pt,itemsep=2pt,parsep=0pt}
\setlist[enumerate]{leftmargin=8mm,topsep=3pt,itemsep=7pt,parsep=0pt}
\pagestyle{fancy}
\fancyhf{}
\fancyhead[L]{\small\color{muted} Дроби и уравнения}
\fancyhead[R]{\small\color{muted} Ксения Васильченко}
\fancyfoot[C]{\small\color{muted}\thepage}
\renewcommand{\headrulewidth}{0.45pt}
\renewcommand{\headrule}{\hbox to\headwidth{\color{accent}\leaders\hrule height\headrulewidth\hfill}}
\newcommand{\sect}[1]{%
  \vspace{5pt}%
  {\Large\color{darkaccent}#1\par}%
  \vspace{2pt}%
  {\color{accent}\rule{\linewidth}{0.55pt}\par}%
  \vspace{5pt}%
}
\newcommand{\minititle}[1]{\par\Needspace{5\baselineskip}\vspace{2pt}{\color{darkaccent}\bfseries #1}\par\vspace{2pt}}
\newcommand{\checkitem}[1]{\item[$\square$] #1}
\newcommand{\important}[1]{{\color{darkaccent}\bfseries #1}}
\begin{document}
\hypersetup{pdftitle={Чек-лист. Дроби и уравнения с дробными коэффициентами},pdfauthor={Лёвин Артём Александрович}}

% Титульный лист
\thispagestyle{empty}
\begin{center}
\vspace*{27mm}
{\Large\color{darkaccent} МАТЕМАТИКА\par}
\vspace{17mm}
{\LARGE\color{textgray} Чек-лист\par}
\vspace{10mm}
{\LARGE\color{accent} Дроби и уравнения\\[3pt] с дробными коэффициентами\par}
\vspace{16mm}
{\normalsize Повторение арифметики и алгебры. 6 класс\par}
\vfill
{\large Ксения Васильченко\par}
\vspace{11mm}
{29 октября 2026 г.\par}
\vspace{18mm}
{\small Лёвин Артём Александрович\\[5pt] эксклюзивно для Ксении Васильченко\par}
\vspace{21mm}
\begin{tikzpicture}[x=1cm,y=1cm]
\draw[accent,line width=.9pt] (-4.1,0) .. controls (-2.3,.65) and (-1.1,-.48) .. (0,0)
 .. controls (1.5,.48) and (2.7,-.65) .. (4.1,0);
\end{tikzpicture}
\vspace{8mm}
\end{center}
\clearpage

\sect{1. Что нужно уметь}
Отметьте пункты, которые можете выполнить самостоятельно.
\begin{itemize}[itemsep=4pt]
\checkitem{Переводить десятичные и смешанные числа в обыкновенные дроби.}
\checkitem{Складывать и вычитать дроби с разными знаменателями.}
\checkitem{Умножать и делить дроби, заранее выполняя сокращения.}
\checkitem{Соблюдать порядок действий, включая скобки и степени.}
\checkitem{Находить неизвестный множитель, делимое и делитель.}
\checkitem{Решать уравнения в несколько действий и приводить подобные слагаемые.}
\end{itemize}

\sect{2. Два вида записи дроби}
\important{Главная идея:} перед вычислением удобно привести дроби к одному виду. Выбирайте запись, в которой числа проще складывать, вычитать, умножать или делить.
\[
0,5=\frac{5}{10}=\frac12,\qquad
0,75=\frac{75}{100}=\frac34,\qquad
1\frac15=\frac{1\cdot5+1}{5}=\frac65.
\]
Обратно, если знаменатель удаётся превратить в $10$, $100$ или $1000$, дробь легко записать десятичной:
\[
\frac9{25}=\frac{9\cdot4}{25\cdot4}=\frac{36}{100}=0,36.
\]
\minititle{Полезные равенства}
\[
\frac12=0,5,\quad \frac14=0,25,\quad \frac34=0,75,\quad
\frac15=0,2,\quad \frac18=0,125.
\]

\sect{3. Действия с обыкновенными дробями}
\minititle{Сложение и вычитание}
Приведите дроби к общему знаменателю, затем сложите или вычтите числители:
\[
\frac16+\frac12+\frac18=\frac4{24}+\frac{12}{24}+\frac3{24}=\frac{19}{24}.
\]
\minititle{Умножение и деление}
Для умножения перемножают числители и знаменатели. При делении на ненулевую дробь умножают на обратную:
\[
\frac{33}{24}:\frac{43}{8}=\frac{33}{24}\cdot\frac8{43}=\frac{11}{43}.
\]
\important{Запомните:} при умножении и делении дробей общий знаменатель искать не требуется.

\sect{4. Порядок действий и скобки}
Сначала выполните действия в скобках и возведение в степень, затем умножение и деление, после этого --- сложение и вычитание.

\minititle{Пример 1. Разные виды дробей в одном выражении}
\[
\begin{aligned}
\left(\frac16+0,5+\frac18\right):3\frac16
 &=\left(\frac16+\frac12+\frac18\right):\frac{19}{6}\\[2pt]
 &=\frac{19}{24}\cdot\frac6{19}=\boxed{\frac14=0,25}.
\end{aligned}
\]
\minititle{Пример 2. Скобки, квадрат и десятичные дроби}
\[
\begin{aligned}
\frac{\left(1\frac15\right)^2-1,08}{0,03}
 &=\frac{1,44-1,08}{0,03}\\[2pt]
 &=\frac{0,36}{0,03}=\frac{36}{3}=\boxed{12}.
\end{aligned}
\]
Здесь $\left(1\frac15\right)^2=1,2^2=1,44$. Чтобы разделить $0,36$ на $0,03$, умножаем делимое и делитель на $100$.

\clearpage
\sect{5. Как найти неизвестное число}
Вспомните связь компонентов арифметических действий. Все деления ниже допустимы только при ненулевом делителе.
\renewcommand{\arraystretch}{1.3}
\begin{center}
\small
\begin{tabularx}{\linewidth}{@{}>{\raggedright\arraybackslash}X>{\raggedright\arraybackslash}X@{}}
\toprule
\textbf{Уравнение} & \textbf{Как найти неизвестное}\\
\midrule
$a\cdot x=c$ & $x=c:a$, если $a\ne0$\\
$a:x=c$ & $x=a:c$, если $a\ne0$, $c\ne0$\\
$x+b=c$ & $x=c-b$\\
$a-x=c$ & $x=a-c$\\
\bottomrule
\end{tabularx}
\end{center}

\minititle{Пример 3. Неизвестный множитель}
\[
\frac27x=1\frac17\quad\Longrightarrow\quad
x=\frac87:\frac27=\frac87\cdot\frac72=\boxed4.
\]
\minititle{Пример 4. Неизвестный делитель}
\[
2:x=1\frac14\quad\Longrightarrow\quad
x=2:\frac54=2\cdot\frac45=\frac85=\boxed{1,6}.
\]
\important{Проверка:} подставьте найденное число в исходное уравнение и убедитесь, что обе части равны.

\sect{6. Уравнения в два действия}
Если рядом с неизвестным есть несколько действий, постепенно освобождайте его от остальных чисел. Сначала уберите прибавленное или вычитаемое число, затем найдите неизвестный множитель.
\minititle{Пример 5}
\[
\begin{aligned}
\frac49x+\frac37&=1,\\
\frac49x&=1-\frac37=\frac47,\\
x&=\frac47:\frac49=\frac47\cdot\frac94=\boxed{\frac97}.
\end{aligned}
\]
Проверка: $\dfrac49\cdot\dfrac97+\dfrac37=\dfrac47+\dfrac37=1$.

\sect{7. Подобные слагаемые}
\important{Подобные слагаемые} имеют одинаковую буквенную часть. Их коэффициенты можно сложить, сохранив общую букву.
\[
\frac19x+\frac49x=\left(\frac19+\frac49\right)x=\frac59x.
\]
\minititle{Пример 6}
\[
\begin{aligned}
\frac19x+\frac49x&=3\frac1{18},\\
\frac59x&=\frac{55}{18},\\
x&=\frac{55}{18}:\frac59=\boxed{\frac{11}{2}=5,5}.
\end{aligned}
\]
\important{Важно:} $\frac19x+\frac49x$ --- это $\frac59x$, а не $\frac59x^2$.

\sect{8. Типичные ошибки}
\begin{itemize}
\item \textbf{Искать общий знаменатель при делении.} Здесь достаточно умножить на обратную дробь.
\item \textbf{Делить на смешанное число без перевода.} Сначала запишите его неправильной дробью.
\item \textbf{Забывать скобки.} Квадрат $\left(1\frac15\right)^2$ относится ко всему смешанному числу.
\item \textbf{Менять только один множитель при делении десятичных дробей.} Умножайте делимое и делитель на одно и то же ненулевое число.
\item \textbf{Менять букву при сложении подобных слагаемых.} Складываются коэффициенты, буквенная часть сохраняется.
\end{itemize}

% Отдельная страница алгоритма
\clearpage
\thispagestyle{fancy}
\begin{center}
{\Large\color{darkaccent}Алгоритм: как решать уравнение с дробями\par}
\vspace{13mm}
\begin{tikzpicture}[
 node distance=10mm,
 >={Stealth[length=2.5mm]},
 every node/.style={align=center,font=\small},
 start/.style={ellipse,draw=accent,thick,minimum width=50mm,minimum height=11mm},
 actionbox/.style={rounded corners=1.5mm,draw=accent,thick,text width=76mm,minimum height=12mm,inner sep=3mm},
 decision/.style={diamond,aspect=2.9,draw=accent,thick,text width=47mm,inner sep=1.8mm},
 arrow/.style={->,draw=darkaccent,thick}
]
\node[start] (start) {Уравнение с дробями};
\node[actionbox,below=of start] (calc) {Выполните числовые действия,\\ если они есть};
\node[actionbox,below=of calc] (like) {Приведите подобные слагаемые\\ с неизвестным, если они есть};
\node[decision,below=12mm of like] (free) {Есть свободное\\ слагаемое?};
\node[actionbox,below=17mm of free] (remove) {Уберите его обратным действием};
\node[actionbox,below=14mm of remove] (isolate) {Найдите неизвестный множитель\\ или делитель};
\node[actionbox,below=of isolate] (check) {Проверьте ответ подстановкой\\ в исходное уравнение};
\node[start,below=of check] (end) {Ответ};
\draw[arrow] (start)--(calc);
\draw[arrow] (calc)--(like);
\draw[arrow] (like)--(free);
\draw[arrow] (free)--node[midway,right,font=\footnotesize] {да} (remove);
\draw[arrow] (free.east) -- ++(36mm,0) |- node[near start,above,font=\footnotesize] {нет} (isolate.east);
\draw[arrow] (remove)--(isolate);
\draw[arrow] (isolate)--(check);
\draw[arrow] (check)--(end);
\end{tikzpicture}
\end{center}

\clearpage
\sect{Тренировочный блок --- Путь А}
\textcolor{muted}{10 заданий. От простого к сложному. Записывайте основные преобразования и сокращайте дроби.}
\vspace{2mm}
\begin{enumerate}[itemsep=11pt]
\item Переведите $0,75$ в обыкновенную дробь, а $\dfrac38$ --- в десятичную.
\item Вычислите: $\displaystyle \frac18+0,5+\frac14$.
\item Выполните деление: $\displaystyle 1\frac23:\frac59$.
\item Вычислите: $\displaystyle \left(\frac16+0,5+\frac13\right):1\frac14$.
\item Вычислите: $\displaystyle \frac{\left(1\frac15\right)^2-1,24}{0,04}$.
\item Решите уравнение: $\displaystyle \frac35x=1\frac45$.
\item Решите уравнение: $\displaystyle 3:x=1\frac15$.
\item Решите уравнение: $\displaystyle \frac23y=1,5-\frac16$.
\item Решите уравнение: $\displaystyle \frac38x+\frac18=1$.
\item Решите уравнение: $\displaystyle \frac29x+\frac56x=3\frac16$.
\end{enumerate}

\clearpage
\sect{Ответы к тренировочному блоку}
\renewcommand{\arraystretch}{1.37}
\begin{center}
\begin{tabular}{@{}c@{\hspace{27mm}}l@{}}
\toprule
\textbf{№} & \textbf{Ответ}\\
\midrule
1 & $\dfrac34$; $0,375$\\
2 & $\dfrac78$\\
3 & $3$\\
4 & $\dfrac45=0,8$\\
5 & $5$\\
6 & $x=3$\\
7 & $x=\dfrac52=2,5$\\
8 & $y=2$\\
9 & $x=\dfrac73=2\dfrac13$\\
10 & $x=3$\\
\bottomrule
\end{tabular}
\end{center}
\vspace{10mm}
\minititle{Краткая памятка}
\begin{tabularx}{\linewidth}{@{}>{\raggedright\arraybackslash}p{55mm}X@{}}
\toprule
\textbf{Что встретилось} & \textbf{Что сделать}\\
\midrule
Сложение/вычитание дробей & Привести к общему знаменателю.\\
Деление дробей & Умножить на обратную ненулевую дробь.\\
Десятичные дроби в частном & Умножить делимое и делитель на $10$, $100$, $1000$ при необходимости.\\
Смешанное число & Перевести в неправильную дробь.\\
Уравнение в несколько действий & Находить неизвестное обратными действиями.\\
Подобные слагаемые & Сложить коэффициенты, сохранить букву.\\
\bottomrule
\end{tabularx}

\vfill
\begin{center}
{\small\color{muted}Лёвин Артём Александрович --- эксклюзивно для Ксении Васильченко}
\end{center}
\end{document}
